\chapter{Sprints 3 à 5~: Demandes RH, Kafka et Communication Temps Réel}

\noindent\textbf{Plan du chapitre}
\begin{enumerate}[nosep]
  \item Sprint 3~: Workflows de demandes RH
  \item Sprint 4~: Prêts, commission de crédit et Apache Kafka
  \item Sprint 5~: Messagerie WebSocket et assistant IA Groq
\end{enumerate}

% ─────────────────────────────────────────────────────────────
\section{Introduction}
% ─────────────────────────────────────────────────────────────

Trois sprints constituent le c\oe{}ur fonctionnel de Synapse. Le premier couvre
les cinq types de demandes RH ordinaires~: congés, formations, documents
administratifs, autorisations. Le deuxième ajoute les demandes financières avec
leur circuit de validation spécifique à la commission de prêt, et introduit Kafka
comme bus d'événements découplant les producteurs des consommateurs de
notifications. Le troisième sprint livre la communication bidirectionnelle par
WebSocket et l'assistant IA conversationnel.

% ─────────────────────────────────────────────────────────────
\section{Sprint 3~: Workflows de demandes RH}
% ─────────────────────────────────────────────────────────────

\subsection{Objectifs}

\begin{table}[H]
\centering
\caption{Périmètre fonctionnel du sprint 3}
\label{tab:s3}
\begin{tabularx}{\textwidth}{|l|L|}
\hline
\rowcolor{headerblue!25}\textbf{Type de demande} & \textbf{Spécificités d'implémentation} \\\hline
\rowcolor{rowgray}
Congé (12 sous-types) & Calcul des jours ouvrés, vérification du solde,
  alerte si dépassement, justificatif médical obligatoire pour \texttt{MALADIE} \\\hline
Formation & Titre, organisme, mode (PRESENTIEL/DISTANCIEL/EN\_LIGNE), budget,
  durée en jours, date planifiée \\\hline
\rowcolor{rowgray}
Document administratif & 8 types (attestation travail, bulletin de paie, CNSS…),
  sélection du nombre de copies et de la langue \\\hline
Autorisation d'absence & Horaires de début et de fin, durée calculée en heures
  par le frontend \\\hline
\rowcolor{rowgray}
Circuit de validation & \texttt{EN\_ATTENTE} → \texttt{VALIDEE\_CHEF} → \texttt{VALIDEE\_RH}~;
  rejet à tout niveau → \texttt{REJETEE} avec commentaire obligatoire \\\hline
\end{tabularx}
\end{table}

\subsection{Diagramme de cas d'utilisation~: dépôt et validation}

\begin{figure}[H]
\centering
\begin{tikzpicture}

  % ─ Acteurs ────────────────────────────────────────────────
  \foreach \y/\name/\lbl in {7.5/emp/Employé, 4.5/chef/Chef, 1.5/rh/Admin RH}{
    \draw[thick] (-2,\y+0.28) circle(0.28cm);
    \draw[thick] (-2,\y) -- (-2,\y-0.7) coordinate(w\name);
    \draw[thick] (w\name)++(0,0.1)--++(-0.3,-0.3) (w\name)++(0,0.1)--++(0.3,-0.3);
    \draw[thick] (w\name)--++(-0.2,-0.45) (w\name)--++(0.2,-0.45);
    \node[font=\small] at (-2,\y-1.4) {\lbl};
  }

  % ─ Frontière ──────────────────────────────────────────────
  \begin{scope}[on background layer]
    \filldraw[umlborder] (0,-0.2) rectangle (10.5,9.5);
    \node[font=\bfseries\small] at (5.25,9.2) {demandes-service};
  \end{scope}

  % ─ UCs ────────────────────────────────────────────────────
  \draw (5,8.5) ellipse(2.3cm and 0.45cm); \node[font=\scriptsize,align=center] at (5,8.5) {Déposer une demande RH};
  \draw (5,7.1) ellipse(2.5cm and 0.45cm); \node[font=\scriptsize,align=center] at (5,7.1) {Sélectionner le type};
  \draw (5,5.8) ellipse(2.5cm and 0.45cm); \node[font=\scriptsize,align=center] at (5,5.8) {Joindre un justificatif};
  \draw (5,4.5) ellipse(2.5cm and 0.45cm); \node[font=\scriptsize,align=center] at (5,4.5) {Annuler (si EN\_ATTENTE)};
  \draw (5,3.2) ellipse(2.5cm and 0.45cm); \node[font=\scriptsize,align=center] at (5,3.2) {Valider → VALIDEE\_CHEF};
  \draw (5,1.9) ellipse(2.5cm and 0.45cm); \node[font=\scriptsize,align=center] at (5,1.9) {Valider finale → VALIDEE\_RH};
  \draw (5,0.7) ellipse(2.5cm and 0.45cm); \node[font=\scriptsize,align=center] at (5,0.7) {Rejeter + commentaire};

  % ─ Relations internes ─────────────────────────────────────
  \draw[->,dashed,>=stealth] (5,8.05)--(5,7.55) node[midway,right,font=\scriptsize]{<<include>>};
  \draw[->,dashed,>=stealth] (5,5.35)--(5,8.5)  node[pos=0.4,right,font=\scriptsize]{<<extend>> [MALADIE]};

  % ─ Associations acteurs ───────────────────────────────────
  \draw[->] (-1.65,7.78) -- (2.7,8.5);
  \draw[->] (-1.65,7.78) -- (2.5,4.5);
  \draw[->] (-1.65,4.78) -- (2.5,3.2);
  \draw[->] (-1.65,4.78) -- (2.5,0.7);
  \draw[->] (-1.65,1.78) -- (2.5,1.9);
  \draw[->] (-1.65,1.78) -- (2.5,0.7);

\end{tikzpicture}
\caption{Diagramme de cas d'utilisation raffiné~: dépôt et validation d'une demande RH}
\label{fig:uc_dem}
\end{figure}

\subsection{Machine à états~: workflow de validation}

\begin{figure}[H]
\centering
\begin{tikzpicture}[node distance=1.4cm]

  \node[umlstate_init] (init) at (0,6) {};
  \node[umlstate] (att)  at (3,6)   {\texttt{EN\_ATTENTE}};
  \node[umlstate] (chef) at (7.5,6) {\texttt{VALIDEE\_CHEF}};
  \node[umlstate, fill=green!15] (rh) at (12,6) {\texttt{VALIDEE\_RH}};
  \node[umlstate, fill=red!15]  (rej) at (7.5,3.5) {\texttt{REJETEE}};
  \node[umlstate, fill=gray!15] (ann) at (3,3.5)   {\texttt{ANNULEE}};

  \draw[trans] (init) -- (att) node[midway,above]{\scriptsize dépôt};
  \draw[trans] (att)  -- (chef) node[midway,above]{\scriptsize Chef approuve};
  \draw[trans] (chef) -- (rh)   node[midway,above]{\scriptsize RH approuve};
  \draw[trans] (att.south)  -- (ann)  node[midway,right]{\scriptsize annulation};
  \draw[trans] (att.south)  to[bend right=10] node[pos=0.4,right]{\scriptsize rejet chef}(rej);
  \draw[trans] (chef.south) -- (rej)  node[midway,right]{\scriptsize rejet RH};
  \draw[trans] (rh)  to[bend left=20] (rej) node[pos=0.4,above]{\scriptsize rejet RH tardif};

  \node[umlstate_final] at (15.5,6) {};
  \draw[trans] (rh) -- (15.2,6);

\end{tikzpicture}
\caption{Machine à états~: cycle de vie d'une demande RH}
\label{fig:fsm_dem}
\end{figure}

\subsection{Réalisation~: formulaire réactif multi-type Angular}

Le composant \texttt{DeposerDemandeComponent} gère un unique
\texttt{FormGroup} dont les validateurs sont activés et désactivés
dynamiquement selon le type sélectionné~:

\begin{lstlisting}[language=Java,caption={Activation conditionnelle des validateurs}]
this.form.get('type')!.valueChanges.subscribe(type => {
  const needsCredit = type === 'CREDIT';
  if (needsCredit) {
    montant.setValidators([Validators.required, Validators.min(1)]);
    dureeMois.setValidators([Validators.required, Validators.max(120)]);
    rembours.setValidators([Validators.requiredTrue]);
    description.clearValidators();
  } else {
    montant.clearValidators(); // ...
  }
  [montant, dureeMois, rembours].forEach(c =>
    c.updateValueAndValidity({ emitEvent: false }));
  this.cdr.markForCheck();
});
\end{lstlisting}

\subsection{Résultats du sprint 3}

\begin{figure}[H]
  \centering
  % \includegraphics[width=\textwidth]{images/sprint3_formulaire.png}
  \caption{Formulaire de dépôt d'une demande RH~: champs conditionnels selon le type sélectionné}
  \label{fig:s3_form}
\end{figure}

% ─────────────────────────────────────────────────────────────
\section{Sprint 4~: Prêts, commission de crédit et Apache Kafka}
% ─────────────────────────────────────────────────────────────

\subsection{Objectifs}

\begin{table}[H]
\centering
\caption{Périmètre du sprint 4}
\label{tab:s4}
\begin{tabularx}{\textwidth}{|l|L|}
\hline
\rowcolor{headerblue!25}\textbf{Tâche} & \textbf{Résultat attendu} \\\hline
\rowcolor{rowgray}
Demandes de prêt & Formulaire avec calcul temps réel de la mensualité et
  de la capacité d'endettement à 33\,\% du salaire \\\hline
Routage conditionnel & Si \texttt{NEEDS\_COMMISSION = 1} et pas de manager~:
  statut initial \texttt{EN\_ETUDE\_DG}~; sinon \texttt{EN\_ATTENTE} \\\hline
\rowcolor{rowgray}
Commission de prêt & Interface de vote FAVORABLE/DÉFAVORABLE~;
  auto-décision au seuil de 2 votes convergents \\\hline
Direction RH & Validation finale des prêts à statut \texttt{VALIDEE\_DG} → \texttt{APPROUVE} \\\hline
\rowcolor{rowgray}
Apache Kafka & Chaque transition de statut publie un \texttt{NotificationEvent}
  sur le topic \texttt{notifications} \\\hline
\end{tabularx}
\end{table}

\subsection{Machine à états~: workflow de prêt}

\begin{figure}[H]
\centering
\begin{tikzpicture}[node distance=1.2cm]

  \node[umlstate_init] (init) at (0,4) {};
  \node[umlstate]            (att)   at (2.8,5.5) {\texttt{EN\_ATTENTE}};
  \node[umlstate]            (etude) at (6.5,4)   {\texttt{EN\_ETUDE\_DG}};
  \node[umlstate]            (vdg)   at (10,5.5)  {\texttt{VALIDEE\_DG}};
  \node[umlstate,fill=green!15](app)  at (13,4)   {\texttt{APPROUVE}};
  \node[umlstate,fill=red!15] (ref)   at (7.5,1.5){\texttt{REFUSE}};

  \draw[trans] (init) -- (att)   node[midway,above,font=\scriptsize]{avec hiérarchie};
  \draw[trans] (init) -- (etude) node[midway,below,font=\scriptsize]{sans hiérarchie};
  \draw[trans] (att)  -- (etude) node[midway,above,font=\scriptsize]{chef ou RH};
  \draw[trans] (etude)-- (vdg)   node[midway,above,font=\scriptsize]{2 votes FAVORABLE};
  \draw[trans] (vdg)  -- (app)   node[midway,above,font=\scriptsize]{Direction RH};
  \draw[trans] (etude.south) -- (ref) node[midway,right,font=\scriptsize]{2 votes DÉFAVORABLE};
  \draw[trans] (vdg.south)   to[bend right=10] (ref) node[midway,right,font=\scriptsize]{rejet RH};

  \node[umlstate_final] at (16,4) {};
  \draw[trans] (app) -- (15.7,4);

\end{tikzpicture}
\caption{Machine à états~: workflow de validation d'une demande de prêt}
\label{fig:fsm_pret}
\end{figure}

\subsection{Architecture Kafka~: notifications asynchrones}

\begin{figure}[H]
\centering
\begin{tikzpicture}[node distance=1.2cm]

  \node[component] (dem) at (0,2) {demandes-service\\:8085};
  \node[infra]     (kaf) at (5,2) {Apache Kafka\\topic: notifications};
  \node[component] (not) at (10,3.2) {notification-service\\:8084};
  \node[component] (ana) at (10,0.8) {analytics-service\\:8083};
  \node[infra]     (mh)  at (14,3.2) {MailHog / SMTP};
  \node[component] (ws)  at (14,0.8) {Push WebSocket};

  \draw[compArr] (dem) -- node[above,font=\scriptsize]{KafkaTemplate.send()} (kaf);
  \draw[compArr] (kaf) -- node[above,font=\scriptsize]{@KafkaListener} (not);
  \draw[compArr] (kaf) -- node[below,font=\scriptsize]{@KafkaListener} (ana);
  \draw[compArr] (not) -- node[above,font=\scriptsize]{e-mail} (mh);
  \draw[compArr] (not) -- node[below,font=\scriptsize]{push} (ws);
  \draw[compArr,dashed] (ana.west) to[bend left=20]
    node[midway,below,font=\scriptsize]{invalidation cache Caffeine} (dem.south);

\end{tikzpicture}
\caption{Pipeline de notifications asynchrones~: demandes-service → Kafka → consommateurs}
\label{fig:kafka}
\end{figure}

Chaque événement est sérialisé en \texttt{NotificationEvent} contenant
\texttt{employeeId}, \texttt{type} (\texttt{VALIDEE\_CHEF}, \texttt{REJETEE}…),
\texttt{referenceId} et \texttt{actionUrl}. Le \textit{notification-service}
consomme sur le groupe \texttt{notifications-group} et notifie simultanément par
WebSocket (destinataire connecté) et par SMTP via MailHog (toujours).

\subsection{Capacité d'endettement~: calcul côté frontend}

\begin{equation}
  \text{capacité} = \lfloor \text{salaire\_mensuel} \times 0{,}33 \rfloor
  \qquad
  \text{taux resultant} = \frac{\text{mensualité}}{\text{salaire\_mensuel}} \times 100
\end{equation}

Le frontend charge le salaire via \texttt{ProfilEmployeService.getProfil()},
calcule la mensualité en temps réel (\texttt{Math.ceil(montant / durée)}) et
affiche un avertissement visuel si la mensualité dépasse la capacité. Le backend
vérifie par une contrainte stricte que le montant demandé ne dépasse pas
$3 \times \text{salaire mensuel brut}$.

\subsection{Résultats du sprint 4}

\begin{figure}[H]
  \centering
  % \includegraphics[width=0.85\textwidth]{images/sprint4_commission.png}
  \caption{Interface de la commission de prêt~: vote et suggestion de montant/tranches}
  \label{fig:s4_comm}
\end{figure}

% ─────────────────────────────────────────────────────────────
\section{Sprint 5~: Messagerie WebSocket et assistant IA Groq}
% ─────────────────────────────────────────────────────────────

\subsection{Objectifs}

\begin{table}[H]
\centering
\caption{Périmètre du sprint 5}
\label{tab:s5}
\begin{tabularx}{\textwidth}{|l|L|}
\hline
\rowcolor{headerblue!25}\textbf{Tâche} & \textbf{Résultat attendu} \\\hline
\rowcolor{rowgray}
Messagerie STOMP & Envoi/réception temps réel~; historique persisté dans
  \texttt{CHAT\_MESSAGES} \\\hline
Notifications push & Compteur de non-lus~; pop-up dès réception~;
  abonnement par rôle (\texttt{/topic/notifications.chef}, …) \\\hline
\rowcolor{rowgray}
Assistant IA & Réponses aux questions RH courantes via Groq
  (\texttt{llama-3.3-70b-versatile})~; affichage différencié IA/humain \\\hline
Mode simulation & \texttt{mockWs: true} dans \texttt{environment.ts}
  substitue un mock de WebSocket pour le développement sans serveur actif \\\hline
\end{tabularx}
\end{table}

\subsection{Diagramme de séquence~: échange WebSocket STOMP}

\begin{figure}[H]
\centering
\begin{tikzpicture}[xscale=1,yscale=0.93]

  \foreach \x/\lbl in {0/Client A (Angular), 3.5/chat-service, 7/Kafka, 10.5/Client B}{
    \node[lifeline,minimum width=2.5cm] at (\x,8.5) {\lbl};
    \draw[thick,dashed] (\x,8.15) -- (\x,0);
  }

  \draw[seqmsg] (0.05,7.5) -- (3.45,7.5) node[midway,above]{\scriptsize CONNECT /ws/chat (SockJS)};
  \draw[seqret] (3.5,6.8)  -- (0.05,6.8) node[midway,above]{\scriptsize CONNECTED};
  \draw[seqmsg] (0.05,6.1) -- (3.45,6.1) node[midway,above]{\scriptsize SUBSCRIBE /topic/messages.conv42};
  \draw[seqmsg] (0.05,5.4) -- (3.45,5.4) node[midway,above]{\scriptsize SEND /app/chat.send \{text, convId\}};

  % Activation chat-service
  \fill[umlblue!60] (3.4,5.1) rectangle (3.6,2.5);

  \draw[seqmsg] (3.5,4.7)  -- (6.95,4.7) node[midway,above]{\scriptsize NotificationEvent → topic notifications};
  \draw[seqmsg] (3.5,4.0)  -- (10.45,4.0) node[midway,above]{\scriptsize MESSAGE /topic/messages.conv42 \{text, sender\}};
  \draw[seqret] (3.5,3.3)  -- (0.05,3.3)  node[midway,above]{\scriptsize MESSAGE (echo)};

  % Assistant IA
  \node[font=\small\itshape, gray] at (5.25,2.0) {— flux assistant IA —};
  \draw[seqmsg] (0.05,1.5) -- (3.45,1.5) node[midway,above]{\scriptsize SEND /app/ai.message \{question\}};
  \draw[seqret] (3.5,0.8)  -- (0.05,0.8) node[midway,above]{\scriptsize MESSAGE /queue/ai-reply \{answer, model: llama-3.3\}};

\end{tikzpicture}
\caption{Diagramme de séquence~: messagerie STOMP et assistant IA Groq}
\label{fig:seq_chat}
\end{figure}

\subsection{Réalisation~: intégration Groq}

Le \textit{chat-service} construit un prompt système contextuel à chaque question
soumise sur \texttt{/app/ai.message}~:

\begin{lstlisting}[language=Java,caption={Construction du prompt Groq dans ChatService}]
String systemPrompt =
  "Tu es l'assistant RH de Synapse. Réponds uniquement aux questions " +
  "relatives aux congés, demandes, formations et procédures RH. " +
  "Langue: français. Sois concis et précis.";

ChatCompletionRequest req = ChatCompletionRequest.builder()
    .model("llama-3.3-70b-versatile")
    .message("system", systemPrompt)
    .message("user", question)
    .temperature(0.3)
    .build();
\end{lstlisting}

La réponse est diffusée sur \texttt{/queue/ai-reply} de l'utilisateur concerné~;
le frontend l'affiche dans un \textit{bubble} distinct (couleur + icône IA)
pour différencier clairement les messages de l'assistant des messages humains.

\subsection{Résultats du sprint 5}

\begin{figure}[H]
  \centering
  % \includegraphics[width=0.85\textwidth]{images/sprint5_chat.png}
  \caption{Interface de messagerie interne~: fil de discussion et réponses de l'assistant IA}
  \label{fig:s5_chat}
\end{figure}

% ─────────────────────────────────────────────────────────────
\section{Conclusion}
% ─────────────────────────────────────────────────────────────

La deuxième release livre le c\oe{}ur opérationnel de Synapse~: cinq types de
demandes avec circuit hiérarchique traçable, un workflow de prêt paramétrable avec
commission dont le routage est déterminé à la création selon deux indicateurs
configurables, Kafka comme bus asynchrone garantissant qu'aucun événement de
transition n'est perdu, et une couche de communication temps réel couplant messagerie
interne et intelligence artificielle conversationnelle. La release~3 complète la
plateforme par la dimension collaborative (projets/tâches) et décisionnelle
(reporting, attrition).
